% from main.tex
%%%%%%%%%%  scan idx 673  |  folio 660  %%%%%%%%%%
effet, dans (i), la premi\`ere assertion r\'esulte de 6.7 (i) ; la deuxi\`eme, de 6.10 et 1.8 (i), et de 6.2 (appliqu\'e \`a $0\to I\to\mathcal O_P\to\mathcal O_X\to0$, $P=\mathbf P^N_k$) et 1.3. Enfin, (ii) r\'esulte de (i) et 1.13.

\textbf{Remarque 6.12.} Soient $k$ un corps alg\'ebriquement clos, et $X$ un $k$-sch\'ema projectif purement de dimension $r$, qui est sans cycle premier associ\'e immerg\'e. Si l'on admet des \'el\'ements nilpotents dans $\mathcal O_X$ de fa\c{c}on arbitraire, on ne peut plus borner les coefficients dans le polyn\^ome de Hilbert $\chi(\mathcal O_X(n))=\sum_{i=0}^{r}a_i\binom{n+i}{i}$ en termes de $a_r$ seulement.

Par exemple, soit $Z$ une courbe lisse de degr\'e $d$ dans $\mathbf P^3_k$. Pour tout $n\gg0$, il existe une surface lisse $Y$ de degr\'e $n$, qui contient $Z$. Alors $Z$ est un diviseur de Cartier sur $Y$; soit $\mathcal O_{X_n}=\mathcal O_Y/\mathcal O_Y(-2Z)$. Enfin, on a la suite exacte $0\to\mathcal O_X(-Z)\to\mathcal O_X\to\mathcal O_Z\to0$ et la formule $\langle Z^2\rangle_Y=2p_a(Z)-2-(n-4)d$ ; d'o\`u que $a_0=\chi(\mathcal O_{X_n}(-1))\to\infty$ ($n\to\infty$), bien que $a_1=2d$.

\textbf{Corollaire 6.13.} \emph{Sous les conditions de 6.1, consid\'erons une suite exacte}
\[
   \mathcal O_P(-m_1)^{\oplus M_1}\longrightarrow
   \mathcal O_P(-m)^{\oplus M}\xrightarrow{\alpha}F\longrightarrow0
\]
\emph{avec $m_1\ge m\ge0$. Alors pour $i=0,\ldots,N$,}
\[
   |a_i|\le A_{N-i}(0,\ldots,0,M;m_1)+M\binom{m}{N-i}.
\]
En effet, soit $I=\operatorname{Ker}(\alpha)$ ; \'evidemment, $I$ est un $(0,\ldots,0,M)$-faisceau, et $\chi(I(n))=\sum_{i=0}^{N}\left[M\binom{-m-1+N-i}{N-i}-a_i\right]\binom{n+i}{i}$ ; d'o\`u l'assertion gr\^ace \`a 6.7 (i), et la formule $\binom{-m-1+Q}{Q}=(-1)^Q\binom{m}{Q}$.

